\section{Training Objectives Through the BP Lens}
\label{sec:training}

Next-token prediction optimizes cross-entropy loss on the training
distribution. This is equivalent to maximum likelihood estimation of the
implicit factor graph parameters. The bayes-learner experiment confirms
that gradient descent finds approximately correct BP weights from this
objective (val MAE 0.000752).

\subsection*{Alternative Objectives}

The BP framework suggests objectives that would produce better-calibrated
Bayes nets:

\begin{itemize}
\item \textbf{Calibration objectives}: penalize overconfident or
underconfident posteriors directly, not just cross-entropy.

\item \textbf{Grounding objectives}: reward outputs that are verifiable
against a declared knowledge base.

\item \textbf{BP consistency objectives}: penalize violations of the BP
fixed-point equations across layers. If beliefs at layer $l+1$ are
inconsistent with what BP would predict given beliefs at layer $l$,
penalize that divergence.
\end{itemize}

\subsection*{The Uniqueness Theorem as a Training Target}

The companion paper proves a uniqueness result: exact posteriors uniquely
force BP weights. This means there is a well-defined target for training
--- the weights that implement exact BP --- and the question is how well
current training objectives approach it.

\subsection*{Epistemic Status}

\begin{center}
\begin{tabular}{ll}
\toprule
Claim & Status \\
\midrule
Gradient descent finds approximate BP weights & Empirical observation \\
Calibration objectives would improve reliability & Architectural proposal \\
BP consistency training would enforce layer-wise calibration & Architectural proposal \\
Uniqueness theorem gives a well-defined training target & Direct consequence \\
\bottomrule
\end{tabular}
\end{center}